\section{Results}

\subsection{The full source domain separates covariance fidelity from discrimination}
The first reviewer-critical test used authentic EEGMMIDB source data from
S001--S008 (360 original task trials) and expanded the scalar path to the
complete \([0,1]\) domain, including the analytic-shrinkage neighbourhood.
With participant-grouped source validation, the macro-AUC optimum occurred at
\(\lambda=0.8\) (0.67833), whereas held-source covariance predictive NLL was
minimised at \(\lambda=0.075\). Ledoit--Wolf and OAS fitted on the final
source representation selected 0.01567 and 0.01953, respectively. The
low-regularisation region was therefore available rather than excluded by
construction.

The separation is visible over the path rather than only at the selected
values. Source macro AUC rises from 0.55732 at \(\lambda=0.015\) to 0.62032
at 0.075, 0.66251 at 0.5 and 0.67833 at 0.8. Covariance NLL moves in the
opposite direction after its optimum: -2.56323 at 0.075, -1.93560 at 0.5 and
-1.64090 at 0.8. The discriminability and covariance one-standard-error
admissible sets overlap only at 0.075 and 0.10 (Jaccard 0.125). Choosing the
covariance-NLL optimum instead of the discriminability optimum costs 0.05801
source macro-AUC; choosing the discriminability optimum increases covariance
NLL by 0.92233. These are source-only mechanism results and do not establish
an unseen-target gain.

\subsection{Strict Z0 execution is feasible, but two targets do not establish efficacy}
The first official EDF feasibility run held out S009 and S010 while using
S001--S008 as source. All 90 target task cues were retained and no target
data entered geometry or parameter selection. On S009, Ledoit--Wolf,
source-participant nested shrinkage and the source-signed/tail controls
obtained AUC 0.57143, 0.57540 and 0.58135, respectively. On S010,
Ledoit--Wolf obtained 0.79960 while nested shrinkage and the signed/tail
controls obtained 0.77183 and 0.77976. This deliberately small block therefore
demonstrates protocol integrity and target heterogeneity, not efficacy.

\subsection{The BNCI objective separation survives a full-domain stress test}
A first BNCI2014\_001 experiment evaluated all nine participants in outer
LOSO with 288 left/right motor-imagery trials per participant. The initial
development grid began at 0.2; hence its apparent benefit over analytic
shrinkage could not answer the reviewer's search-space objection. After those
target outcomes had been exposed, we therefore ran a disclosed post-hoc
stress audit over the complete \([0,1]\) domain. It is not treated as
independent confirmation.

Across the nine outer folds, analytic Ledoit--Wolf selected
\(\lambda=0.00227\)--0.00355. The source covariance-NLL optimum was 0.10 in
all nine folds, whereas the source participant-macro-AUC optimum was 0.50 in
one fold, 0.90 in two and 0.95 in six. Thus restoring the near-zero region did
not move the discriminability optimum towards analytic shrinkage.
Participant-macro target AUC was 0.75291 for Ledoit--Wolf, 0.75295 for OAS,
0.76475 for logistic regression, 0.77182 for the source covariance-NLL
optimum and 0.78646 for the full-domain source-discriminability optimum.
The latter exceeded Ledoit--Wolf by a descriptive mean 0.03355 AUC (7/9
participants improved; exact sign-flip \(p=0.03125\)). Because the same nine
targets had been exposed previously, this effect is classified as post-hoc
mechanism stress evidence rather than a confirmatory generalisation result.

\subsection{Prospective falsification removes two candidate replacement rules}
The preregistered lower-tail signed selector did not improve over nested
source-participant scalar tuning in the initial BNCI2014\_001 run: its
participant-macro AUC was 0.78455 versus 0.78782 for nested source-macro-AUC
tuning, a mean paired difference of -0.00328, with a worst participant
difference of -0.04070. It is therefore retained as a direction-aware
mechanism control rather than a superior algorithm.

A separate candidate, covariance-fidelity-constrained transfer shrinkage
(CFCTS), was frozen before evaluating EEGMMIDB S011--S020. The source set
remained S001--S008 and the information budget was Z0. Mean AUC was 0.66730
for the preregistered nested comparator and 0.63309 for CFCTS, a paired mean
difference of -0.03421 (participant bootstrap 95\% CI
[-0.05593,-0.01190], exact two-sided sign-flip \(p=0.02734\)). Six of ten
participants were harmed by more than 0.02 AUC and the worst difference was
-0.08893. CFCTS is therefore rejected as the primary correction rather than
retuned on these participants.

\subsection{Prospective full-domain mismatch predicts an untouched 20-participant target block}
The full-domain S021--S040 experiment was frozen before any target outcome.
Source participants remained S001--S008; the target-information budget was
Z0; every original task cue was retained. The source-only path selected
\(\lambda_D=0.8\) for participant discriminability and \(\lambda_C=0.075\)
for covariance predictive NLL, while the final source Ledoit--Wolf and OAS
values were 0.01567 and 0.01953.

Across the 20 untouched target participants, participant-macro AUC was
0.62919 for Ledoit--Wolf, 0.63562 for OAS, 0.67003 for logistic regression,
0.66509 for source covariance-NLL selection and 0.72116 for full-domain
source-discriminability selection. Relative to Ledoit--Wolf, the latter
improved mean participant AUC by 0.09197 (median 0.07133; 16/20 participants
improved; participant-bootstrap 95\% CI [0.04108,0.14310]; exact
\(2^{20}\) sign-flip \(p=0.00255\); Holm-adjusted \(p=0.00510\)).
Relative to the stronger source covariance-NLL comparator, the improvement
was 0.05607 (median 0.05929; 15/20 improved; 95\% CI
[0.02185,0.09049]; exact sign-flip and Holm-adjusted \(p=0.00597\)).

The effect was not universal. Four participants were harmed by more than
0.02 AUC relative to Ledoit--Wolf and five relative to covariance-NLL
selection; the worst paired changes were -0.11067 and -0.07312,
respectively. The prospective result therefore supports a target consequence
of the objective mismatch while motivating a separate hard-participant
reliability analysis.

A stricter post-primary delete-one-source sensitivity propagated source-pool
variation through the entire analysis rather than fixing the original
hyperparameter. After omitting each of S001--S008 in turn, we recomputed every
inner source tangent reference, reran the complete \([0,1]\) path, reselected
both the discriminability and covariance-NLL optima, rebuilt the final tangent
reference and refit the classifiers. The discriminability-selected strength
varied across \(\{0.4,0.8,0.95\}\), whereas the covariance-NLL optimum was
0.075 in seven omissions and 0.10 in one, demonstrating genuine source-pool
selection uncertainty. Nevertheless, discriminability-selected macro AUC was
0.71272 on average across omissions (SD 0.00926; range
0.69999--0.72581). Its advantage over a freshly reselected covariance-NLL
model remained positive in 8/8 omissions (mean +0.05556; range
+0.04329 to +0.07172), and its advantage over a freshly refit Ledoit--Wolf
model was also positive in 8/8 (mean +0.10079; range +0.07069 to +0.13741).
This is a post-primary training-pool sensitivity analysis, not an additional
confirmatory hypothesis test, but it directly shows that the prospective
effect does not depend on one source participant or on freezing the originally
selected lambda.

A separate deterministic score-regeneration audit then rebuilt all five
prospective target methods from the immutable S021--S040 trial-covariance
artifact and the frozen implementation semantics. Across all 20 target
participants and five methods, regenerated AUCs matched the workflow table at
the stored eight-decimal precision (maximum absolute difference 0.0).
This closes a result-lineage gap without rewriting the original chronology:
the original workflow did not archive an independent trial-score file, so the
regenerated score hashes are explicitly post-hoc reproducibility evidence
rather than an original score freeze.

\subsection{Two TRUE PRE-DATA external cohorts replicate source-objective separation and preserve the target-effect direction}
The external v5 panel was frozen before project access to either
BNCI2014\_002 or BNCI2015\_001 outcomes. The scientific runner, source
selection rules, dense path, secondary selectors and score-freeze contract
were fixed before the first cohort was opened. BNCI2015\_001 was then
executed with byte-identical scientific assets; its execution-only child PR
changed only the CI trigger. BNCI2014\_002 contributed 14 participants,
15 native EEG channels and 2240 target trials. BNCI2015\_001 contributed
12 participants, 13 native channels and 5600 target trials across two or
three sessions per participant. Both analyses were Z0.

Source-objective separation was complete across the external panel. In
BNCI2014\_002, source participant-discrimination and covariance-NLL selected
different strengths in 14/14 outer folds: the discrimination optimum lay
between 0.7 and 0.9, while the covariance optimum was zero in 13 folds and
0.02 in one. In BNCI2015\_001 they differed in 12/12 folds: the
discrimination optimum was 0.9 or 0.95, while covariance-NLL selected
0.015--0.05. Thus the two TRUE PRE-DATA cohorts contain 26/26 outer folds
with different source objective optima. This fold count is descriptive
consistency within two datasets, not 26 independent replication strata.

The target consequence had the same direction in both datasets but the
frozen small-sample exact tests remained individually inconclusive. In
BNCI2014\_002, participant-macro AUC was 0.77810 for source-macro selection,
0.76047 for covariance-NLL and 0.75523 for Ledoit--Wolf.
Source-macro minus covariance-NLL was +0.01763 AUC
(95\% participant-bootstrap CI [-0.00358,0.03728], exact sign-flip
\(p=0.12329\), Holm \(p=0.13965\)); source-macro minus Ledoit--Wolf was
+0.02287 (95\% CI [0.00045,0.04402], Holm \(p=0.13965\)).
In BNCI2015\_001, the corresponding AUCs were 0.76697, 0.72368 and 0.72166.
The source-macro advantages were +0.04329 over covariance-NLL
(95\% CI [0.00016,0.09345], exact \(p=0.10107\), Holm
\(p=0.18555\)) and +0.04531 over Ledoit--Wolf
(95\% CI [0.00290,0.09565], exact \(p=0.09277\), Holm
\(p=0.18555\)). The equal-dataset descriptive means of H2 and H3 were
+0.03046 and +0.03409 AUC, respectively. We therefore describe the target
effect as directionally replicated across both TRUE PRE-DATA cohorts, not as
two individually significant superiority confirmations.

The harmed-participant tail persisted. For H2, three of 14 BNCI2014\_002
participants and two of 12 BNCI2015\_001 participants lost more than
0.02 AUC; the corresponding CVaR25 values were negative. The predeclared
mean--tail--fidelity selector therefore remained secondary: its change versus
pure source-macro selection was +0.00416 in BNCI2014\_002 and +0.00055 in
BNCI2015\_001, with negative target CVaR25 in both. A descriptive audit of
the simpler discriminability-first lexicographic selector also failed to
support promotion, with mean changes of -0.00865 and -0.00753 AUC relative
to pure source-macro selection. The external evidence therefore strengthens
the objective-mismatch finding without producing a superior new selector.

BNCI2015\_001 additionally supplied the predeclared descriptive
cross-session check. Source-macro minus covariance-NLL was positive on average
within sessions 0A, 1B and 2C (+0.06252, +0.05296 and +0.06270 AUC,
respectively); source-macro minus Ledoit--Wolf was similarly positive
(+0.06495, +0.05432 and +0.06657). Session rows were not used as independent
inferential units.

\subsection{Post-outcome participant-harm uncertainty and influence}
To check whether positive cohort means conceal negative transfer, we performed
an explicitly post-outcome sensitivity audit of all six frozen participant
effect vectors (three cohorts by two covariance comparators). Relative to
source covariance-NLL, the numbers of participants losing more than 0.02 AUC
were 5/20 in EEGMMIDB S021--S040, 3/14 in BNCI2014\_002 and 2/12 in
BNCI2015\_001; the corresponding Ledoit--Wolf counts were 4/20, 2/14 and
2/12. For the covariance-NLL contrasts, 95\% Wilson iid-participant
reference intervals for those harm fractions were [0.112, 0.469],
[0.076, 0.476] and [0.047, 0.448], respectively. These intervals are
conditional descriptive references, not source-pool or new-dataset safety
guarantees, because the target participants share fitted source models.

The frozen 46-person individual effect distribution is visualised in
Fig.~\ref{fig:participant-harm-v10}; the plot is descriptive rather than
an additional independent experiment.

An additional post-outcome influence calculation removed one target
participant at a time without refitting the shared source model. The smallest
remaining mean effects versus covariance-NLL were +0.0475, +0.0142 and
+0.0258 AUC in the same dataset order. Thus the observed positive
within-cohort mean directions do not depend on any single test participant.
However, BNCI2015\_001 had a mean gain of +0.0433 but a median gain of only
+0.0105, indicating a right-skewed benefit distribution. Neither this
delete-one diagnostic nor its positive direction supplies a new
confirmatory test or resolves the substantial harmed-participant tail.
The exact participant-level audit is reproducible by
src/participant\_harm\_influence\_audit.py using the immutable three-cohort
paired-effect receipt.

\subsection{Corrected TRUE PRE-DATA validity frontiers show reproducible decision variation but cohort-specific steepness}
The original v5 runner emitted a generic zero-crossing statistic for an AUC
path centered on its own covariance optimum. That statistic was identified as
structurally uninformative before the first v5 target result became available:
the covariance-optimal point has zero centered difference by construction and
therefore places zero inside every non-negative admissibility set. Following
the frozen pre-outcome correction, we instead report the complete
covariance-regret--AUC envelope, its maximum span, and the
material-fragility budget
\[
 \delta_\tau=
 \inf\{\delta:
 \operatorname{span}(E_{AUC}(\delta))\ge\tau\},
\]
with \(\tau=0.02\) AUC primary and 0.01/0.05 descriptive. Each
20,000-draw participant bootstrap resamples outer participants once and reuses
the same indices at every \(\lambda\).

BNCI2014\_002 had a maximum dense-path AUC envelope span of 0.03844
(bootstrap 95\% CI [0.02714,0.06656]). Its primary material-fragility budget
was \(\delta_{0.02}=0.91611\) raw per-dimension covariance-NLL regret,
finite in 99.945\% of bootstrap draws; the conditional-finite 95\% interval
was [0.26093,1.29696]. The 0.05 threshold was not reached in the point
estimate. BNCI2015\_001 had a larger maximum span of 0.05011
(95\% CI [0.02238,0.10335]) and reached 0.02 at a substantially smaller
budget, \(\delta_{0.02}=0.17519\), finite in 98.77\% of draws
(conditional-finite 95\% CI [0.04263,0.61685]); its point
\(\delta_{0.05}\) was 0.57390.

The result is robust to the predeclared grid-resolution audit. Reducing the
405-point path to every eighth candidate changes the maximum span only from
0.03844 to 0.03825 in BNCI2014\_002 and from 0.05011 to 0.05004 in
BNCI2015\_001; the corresponding \(\delta_{0.02}\) values become 0.93097
and 0.17788. Thus both cohorts support a non-trivial regularisation validity
frontier, while the markedly different critical budgets reject a universal
frontier steepness or universal fidelity threshold.

\subsection{Simultaneous-path sensitivity limits material-frontier inference}
The corrected validity frontiers have observed maximum AUC spans of 0.0384
and 0.0501 in the two external cohorts. Their paired-participant percentile
bootstrap intervals quantify resampling variability of the maximum-spread
statistic, but they are not simultaneous confidence bands over the 405
candidate shrinkage values. In particular, an empirical maximum--minimum
span can be positive under a flat population AUC path merely because many
noisy candidate points are searched. The percentile interval's lower
endpoint exceeding 0.02 does not itself justify a population-level
materiality claim.

A post-outcome sensitivity implementation now estimates the supremum of the
participant-bootstrap mean-path deviation across all shrinkage candidates.
Writing \(\widehat A(\lambda)\) for the observed participant-mean AUC path
and \(q_{.95}\) for the 95th percentile of
\(\sup_\lambda|\widehat A^*(\lambda)-\widehat A(\lambda)|\), we report the
conservative simultaneous-path sensitivity lower bound
\[
  L_{\mathrm{span}}=
  \max\left\{0,
    \big[\max_\lambda\widehat A(\lambda)-
         \min_\lambda\widehat A(\lambda)\big]-2q_{.95}
  \right\}.
\]
This is a nominal post-outcome diagnostic conditional on the estimated
source covariance-regret axis; it does not represent refitting uncertainty
from overlapping leave-one-person-out source pools. No numerical
simultaneous-bound result is asserted until the archived dense-path AUC
matrices are reprocessed. The descriptive frontier observations remain
valid, but material population-level fragility is a separate inferential
question.

\subsection{Real-artifact V11 reanalysis closes the 405-point inference question}
We recovered both original TRUE PRE-DATA v5 GitHub Actions artifacts,
verified their raw result JSON SHA-256 hashes, and independently reconstructed
the 14-by-405 and 12-by-405 participant-by-regularisation AUC matrices.
The observed range and the original paired percentile bootstrap intervals
reproduced exactly up to numeric precision. BNCI2014\_002 retained the
observed span 0.03844 and originally reported percentile interval
[0.02714,0.06656]; BNCI2015\_001 retained span 0.05011 and interval
[0.02238,0.10335]. Both remain observations, not simultaneously
corrected population effect statements.

We next used 20,000 participant-paired whole-path resamples, with seed
20261010, and evaluated a translation-invariant path error
\(\operatorname{osc}(\bar A^*-\bar A)\), the difference between its maximum
and minimum across all 405 candidates. Its nominal 95th percentile was
0.04035 for BNCI2014\_002 and 0.05923 for BNCI2015\_001.
Because \(|\operatorname{osc}(f)-\operatorname{osc}(g)|
\leq\operatorname{osc}(f-g)\), a conservative conditional span sensitivity
lower bound is the positive part of the observed span minus this radius.
The resulting lower bound is zero in both cohorts; neither supplies
simultaneously adjusted evidence that the underlying participant-mean
path span exceeds 0.02 AUC.

This is a post-outcome conditional bootstrap sensitivity analysis; shared
outer source fits, their model-selection variation, and sampling an
additional dataset are not included. It does not establish that the true
path is flat. It prevents us from using selected extrema as a confirmatory
materiality result. The separately predeclared cross-participant primary
selection contrasts are unchanged by this correction.

\subsection{Prospective representation testing falsifies the historical hidden-filter-bank narrative}
The S041--S060 representation block was frozen independently of S021--S040.
It compared WIDE231, FB8-1848, source-only PCA-compressed FB8-PCA231 and a
label-independent WIDE+NUISANCE1848 dimension control under Ledoit--Wolf,
OAS, full-domain source discrimination and full-domain source covariance-NLL.
All 20 targets remained Z0 and all original task cues were retained in every
band.

Under source-discrimination selection, participant-macro AUC was 0.74461 for
WIDE231, 0.66202 for FB8-1848, 0.68884 for FB8-PCA231 and 0.62002 for
WIDE+NUISANCE1848. The preregistered primary interaction,
\[
 [A_{\mathrm{FB8},D}-A_{\mathrm{WIDE},D}]
 -
 [A_{\mathrm{FB8},LW}-A_{\mathrm{WIDE},LW}],
\]
was -0.07781 AUC (95\% CI [-0.12415,-0.03451], exact sign-flip
\(p=0.00332\), Holm \(p=0.01328\)). The corresponding interaction using
source covariance-NLL as the covariance-oriented comparator was -0.04312
(95\% CI [-0.06794,-0.01741], Holm \(p=0.01751\)). Thus regularisation
objective choice prospectively changes the representation comparison, but in
the opposite direction to the historical hypothesis: discrimination-oriented
regularisation strengthens the WIDE advantage rather than revealing a hidden
filter-bank advantage.

The matched-dimension control also contradicts a pure spectral-information
story. FB8-PCA231 minus WIDE231 under source-discrimination selection was
-0.05577 (95\% CI [-0.09992,-0.01380], Holm \(p=0.04633\)).
Conversely, FB8 minus the label-independent WIDE+NUISANCE1848 control was
+0.04200 but did not survive the four-contrast family (Holm
\(p=0.08661\)). The frozen source-budget intervention supplied a second
fixed-signal counterfactual: with 25\%, 50\%, 75\% and 100\% of source
trials retained, the source-discrimination advantage over covariance-NLL was
+0.04370, +0.08515, +0.06923 and +0.07755 AUC. Neither effect size nor
selected shrinkage strength varied monotonically with source sample budget.
Together, these results prospectively reject both the historical
filter-bank-reversal headline and a simple monotonic \(p/n\) mechanism.

\subsection{A target-calibration AUC experiment exposes an identifiability limitation}
A separate prospective EEGMMIDB block, S061--S080, compared Z0
source-only scoring with unlabeled U5/U10/U20-prefix and U-all target
adaptation. Under Z0, source-discriminability selection improved over
Ledoit--Wolf by +0.07463 macro AUC (20 participants; 95\% bootstrap CI
[0.03695,0.11042]; Holm \(p=0.01296\)), with 17 participants improved
and two harmed by more than 0.02. The nine predeclared deployment-budget
contrasts passed the structural audit.

All matched U-k and U-all changes in ROC AUC were exactly zero, but this
result is \emph{not} evidence that unlabeled calibration is empirically
useless. The frozen calibration applied only a constant target-participant
feature shift \(b_k\) before a fixed linear score \(s(x)=w^\top(x-c)\).
For every trial of that participant,
\[
s_k(x)=w^\top(x-b_k-c)=s_0(x)-w^\top b_k.
\]
The subtraction is independent of the trial and therefore leaves the
within-participant score ordering unchanged. Since ROC AUC depends
only on this ordering, its paired U-k/Z0 contrast is algebraically forced
to zero. These zero effects are retained as an identifiability negative
control, not interpreted as a measured calibration null. Any assessment of
calibration utility must use a source-frozen operating threshold or a
threshold-sensitive utility measure, and an S061--S080 extension would
necessarily be post-outcome.


\subsection{Dataset-stratified effect consistency is descriptive, not a third prospective test}
To inspect replication-stratum heterogeneity without pseudoreplicating
trials, we extracted all 46 held-participant paired effects for the three
prospectively evaluated cohorts (EEGMMIDB S021--S040, BNCI2014\_002,
BNCI2015\_001) from their original successful workflows. The original
primary H2 (source-discriminability minus covariance-NLL) cohort mean
differences were +0.05607, +0.01763 and +0.04329, respectively.
The primary H3 (source-discriminability minus Ledoit--Wolf) differences
were +0.09197, +0.02287 and +0.04531.

A newly designed, explicitly \emph{post-outcome descriptive}
equal-dataset summary gives +0.03900 AUC for H2 and +0.05338 for H3;
restricting it to the two independent external cohorts gives +0.03046
and +0.03409, respectively. These values assign one weight per dataset
regardless of its participant/trial count; they are not participant- or
trial-pooled effects. The companion conditional participant-resampling
audit, with 20,000 seeded draws inside each fixed dataset, yields
intervals [0.01938,0.06022] for the three-dataset H2 descriptive mean,
and [0.02989,0.07813] for H3. Across the two external datasets alone
the conditional intervals are [0.00696,0.05754] and
[0.01002,0.06183].

These conditional intervals do \emph{not} account for uncertainty about
sampling new EEG acquisition datasets, overlapping source fits or the
post-outcome choice to combine cohorts. They cannot supersede the
original cohort-wise Holm results, which did not establish independent
external superiority. Moreover, the source-macro H2 method harmed
10 of the 46 held participants by more than 0.02 AUC relative to
covariance-NLL (5/20, 3/14, 2/12), and H3 harmed 8/46 relative to
Ledoit--Wolf (4/20, 2/14, 2/12). This replication-stratified pattern
supports bounded objective-mismatch conclusions with explicit harm,
not a universally safe new shrinkage rule.
